% Part: normal-modal-logic
% Chapter: frame-definability
% Section: definability

\documentclass[../../../include/open-logic-section]{subfiles}

\begin{document}

\olfileid[pt]{nml}{frd}{def}

\olsection{Definibilidade de frames}

Embora as implicações recíprocas de \olref[acc]{thm:soundschemas} falhem,
elas valem se substituirmos «modelo» por «frame»: para as propriedades
consideradas em \olref[acc]{thm:soundschemas}, \emph{é} verdade que, se
!!a{formula} é válida em um \emph{frame}, então a relação de acessibilidade
desse frame possui a propriedade correspondente. Assim, as !!{formula}s
consideradas \emph{definem} as classes de frames que possuem a propriedade
correspondente.

\begin{defn}
  Se $\mClass{F}$ é uma classe de frames, dizemos que $!A$
  \emph{define}~$\mClass{F}$ se e somente se $\mModel{F} \Entails !A$
  exatamente para os frames~$\mModel{F} \in \mClass{F}$.
\end{defn}

Passamos agora a estabelecer os resultados completos de definibilidade para
frames.

\begin{thm}\ollabel{thm:fullCorrespondence}
Se a !!{formula} no lado direito de \olref[acc]{tab:five} é válida em um
frame~$\mModel{F}$, então $\mModel{F}$ possui a propriedade do lado esquerdo.
\end{thm}

\begin{proof}
  \begin{enumerate}
  \item Suponha que \Ax{D} seja válida em
    $\mModel{F} = \tuple{W, R}$, isto é,
    $\mModel{F} \Entails \Box p \lif \Diamond p$. Seja
    $\mModel{M} = \tuple{W, R, V}$ um modelo baseado em~$\mModel{F}$, e seja
    $w \in W$. Precisamos mostrar que existe um $v$ tal que~$Rwv$. Suponha
    que não: então temos tanto $\mSat{M}{\Box !A}$ como
    $\mSat/{M}{\Diamond !A}[w]$ para qualquer~$!A$, inclusive~$p$. Mas, nesse caso,
    $\mSat/{M}{\Box p \lif \Diamond p}[w]$, contrariando a hipótese de que
    $\mModel{F} \Entails \Box p \lif \Diamond p$.
  \item Suponha que \Ax{T} seja válida em $\mModel{F}$, isto é,
    $\mModel{F} \Entails \Box p \lif p$. Seja $w \in W$ um mundo arbitrário;
    precisamos mostrar que $Rww$. Defina $u \in V(p)$ se e somente se $Rwu$
    (quando $q$ é distinto de $p$, $V(q)$ é arbitrário, por exemplo,
    $V(q) = \emptyset)$. Seja $\mModel{M} = \tuple{W, R, V}$. Por construção,
    para todo $u$ tal que $Rwu$: $\mSat{M}{p}[u]$ e, portanto,
    $\mSat{M}{\Box p}[w]$. Mas, por hipótese, $\Box p \lif p$ é verdadeira
    em~$w$, de modo que $\mSat{M}{p}[w]$; pela definição de $V$, porém, isso
    só é possível se~$Rww$.
  \item Demonstramos a contrapositiva: suponha que $\mModel{F}$ não seja
    simétrico; mostraremos que \Ax{B}, isto é, $p \lif \Box\Diamond p$, não é
    válida em $\mModel{F}= \tuple{W, R}$. Se $\mModel{F}$ não é simétrico,
    existem $u$, $v \in W$ tais que $Ruv$, mas não $Rvu$. Defina $V$ de modo
    que $w \in V(p)$ se e somente se não $Rvw$ (e que $V$ seja arbitrária nos
    demais casos). Seja $\mModel{M} =\tuple{W, R, V}$. Agora, pela definição de
    $V$, $\mSat{M}{p}[w]$ para todo $w$ tal que não $Rvw$; em particular,
    $\mSat{M}{p}[u]$, pois não~$Rvu$. Além disso, como $Rvw$ se e somente se
    $w \notin V(p)$, não existe~$w$ tal que $Rvw$ e $\mSat{M}{p}[w]$ e,
    portanto, $\mSat/{M}{\Diamond p}[v]$. Como~$Ruv$, também
    $\mSat/{M}{\Box\Diamond p}[u]$. Segue-se que
    $\mSat/{M}{p \lif \Box\Diamond p}[u]$ e, assim, $\Ax{B}$ não é válida
    em~$\mModel{F}$.
  \item Suponha que \Ax{4} seja válida em $\mModel{F} = \tuple{W,R}$, isto é,
    $\mModel{F} \Entails \Box p \lif \Box\Box p$, e sejam $u$, $v$, $w \in W$
    mundos arbitrários tais que $Ruv$ e $Rvw$; precisamos mostrar que $Ruw$.
    Defina $V$ de modo que $z \in V(p)$ se e somente se $Ruz$ (e que $V$ seja
    arbitrária nos demais casos). Seja $\mModel{M} =\tuple{W, R, V}$. Pela
    definição de $V$, $\mSat{M}{p}[z]$ para todo $z$ tal que $Ruz$ e,
    portanto, $\mSat{M}{\Box p}[u]$. Mas, pela hipótese \Ax{4},
    $\Box p \lif \Box \Box p$ é verdadeira em $u$, de modo que
    $\mSat{M}{\Box \Box p}[u]$. Como $Ruv$ e $Rvw$, temos
    $\mSat{M}{p}[w]$; pela definição de~$V$, porém, isso só é possível se
    $Ruw$, como desejado.
  \item Procedemos contrapositivamente, supondo que o frame
    $\mModel{F} = \tuple{W, R}$ não seja euclidiano, e mostramos que ele
    falsifica~\Ax{5}, isto é,
    $\mModel{F} \Entails/ \Diamond p \lif \Box\Diamond p$. Suponha que existam
    mundos $u$, $v$, $w \in W$ tais que $Rwu$ e $Rwv$, mas não~$Ruv$. Defina
    $V$ de modo que, para todo mundo $z$, $z \in V(p)$ se e somente se
    \emph{não} ocorre~$Ruz$. Seja $\mModel{M} =\tuple{W, R, V}$. Então, por
    hipótese, $\mSat{M}{p}[v]$ e, como $Rwv$, também
    $\mSat{M}{\Diamond p}[w]$. Contudo, não existe mundo $y$ tal que $Ruy$ e
    $\mSat{M}{p}[y]$, de modo que $\mSat/{M}{\Diamond p}[u]$. Como $Rwu$,
    segue-se que $\mSat/{M}{\Box\Diamond p}[w]$, e, portanto, \Ax{5},
    $\Diamond p \lif \Box\Diamond p$, falha em~$w$.
  \end{enumerate}
\end{proof}

Você perceberá uma diferença entre a demonstração para \Ax{D} e os demais
casos: não se fez menção à valoração~$V$. Com efeito, demonstramos que, se
$\mSat{M}{\Ax{D}}$, então $\mModel{M}$ é serial. Assim, \Ax{D} define a classe
dos \emph{modelos} seriais, não apenas dos frames.

\begin{cor}\ollabel{prop:D-serial}
  Todo modelo no qual \Ax{D} é verdadeira é serial.
\end{cor}

\begin{cor}
Cada !!{formula} no lado direito de \olref[acc]{tab:five} define a classe de
frames que possuem a propriedade do lado esquerdo.
\end{cor}

\begin{proof}
  Em \olref[acc]{thm:soundschemas}, demonstramos que, se um modelo possui a
  propriedade da esquerda, a !!{formula} da direita é verdadeira nele.
  Portanto, se um frame~$\mModel{F}$ possui a propriedade da esquerda, a
  !!{formula} da direita é válida em~$\mModel{F}$. Em
  \olref{thm:fullCorrespondence}, demonstramos as implicações recíprocas: se
  !!a{formula} da direita é válida em~$\mModel{F}$, então $\mModel{F}$ possui
  a propriedade da esquerda.
\end{proof}

\begin{prob}
Mostre que, se a !!{formula} no lado direito de
\olref[nml][frd][acc]{tab:anotherfive} é válida em um frame~$\mModel{F}$,
então $\mModel{F}$ possui a propriedade do lado esquerdo. Para isso, considere
um frame que \emph{não} satisfaça a propriedade da esquerda e defina uma~$V$
adequada, de modo que a !!{formula} da direita seja falsa em algum mundo.
\end{prob}

\olref{thm:fullCorrespondence} também mostra que as propriedades podem ser
combinadas: por exemplo, se tanto \Ax{B} como \Ax{4} são válidas
em~$\mModel{F}$, então o frame é simétrico e transitivo etc. Muitas lógicas
modais importantes são caracterizadas como o conjunto de !!{formula}s válidas
em todos os frames que combinam certas propriedades de frames; portanto,
podemos caracterizá-las como o conjunto de !!{formula}s válidas em todos os
frames nos quais são válidas as !!{formula}s definidoras correspondentes. Por
exemplo, o sistema clássico \Log{S4} é o conjunto de todas as !!{formula}s
válidas em todos os frames reflexivos e transitivos, isto é, em todos aqueles
nos quais tanto \Ax{T} como~\Ax{4} são válidas. \Log{S5} é o conjunto de todas
as fórmulas válidas em todos os frames reflexivos, simétricos e euclidianos,
isto é, em todos aqueles nos quais \Ax{T}, \Ax{B} e \Ax{5} são válidas.

Relações lógicas entre propriedades de $R$ em geral correspondem a relações
entre as !!{formula}s definidoras correspondentes. Por exemplo, toda relação
reflexiva é serial; portanto, sempre que \Ax{T} é válida em um frame, \Ax{D}
também é. (Observe que essa relação \emph{não} é uma relação de consequência.
Não ocorre que, sempre que $\mSat{M}{\Ax{T}}[w]$, então
$\mSat{M}{\Ax{D}}[w]$.) Registramos algumas relações desse tipo.

\begin{prop}\ollabel{prop:relation-facts}
  Seja $R$ uma relação binária sobre um conjunto $W$; então:
  \begin{enumerate}
  \item Se $R$ é reflexiva, então ela é serial.
  \item Se $R$ é simétrica, então ela é transitiva se e somente se é
    euclidiana.
  \item Se $R$ é simétrica ou euclidiana, então ela é fracamente dirigida
    (possui a «propriedade do losango»).
  \item Se $R$ é euclidiana, então ela é fracamente conexa.
  \item Se $R$ é funcional, então ela é serial.
  \end{enumerate}
\end{prop}

\begin{prob}
  Demonstre \olref[nml][frd][def]{prop:relation-facts}.
\end{prob}

\end{document}
